% Generated draft: requires mathematical and source review.
\documentclass[11pt]{article}
\usepackage[margin=25mm]{geometry}
\usepackage{fontspec}
\setmainfont{Latin Modern Roman}
\usepackage{amsmath,amssymb,mathtools,unicode-math}
\setmathfont{Latin Modern Math}
\usepackage{enumitem,needspace,xurl,hyperref}
\setlength{\parindent}{0pt}
\setlength{\parskip}{5pt}
\setlength{\emergencystretch}{3em}
\newcommand{\N}{\mathbb{N}}
\newcommand{\Z}{\mathbb{Z}}
\newcommand{\Q}{\mathbb{Q}}
\newcommand{\R}{\mathbb{R}}
\newcommand{\C}{\mathbb{C}}
\newcommand{\F}{\mathbb{F}}
\newcommand{\A}{\mathbb{A}}
\newcommand{\PP}{\mathbb{P}}
\newcommand{\E}{\mathbb{E}}
\newcommand{\eps}{\varepsilon}
\newcommand{\dd}{\,\mathrm{d}}
\begin{document}
\section*{UnsolvedMath ID: 10900028}
\textbf{Source problem code:} catalogue entry
\subsection{Definitions and mathematical statement}
For the usual orientable setting, let $M$ be a complete finite-volume hyperbolic $3$-manifold and let $N$ be a compact core, with torus boundary in the cusped case and $N=M$ in the closed case. For an integral class $a\in H_2(N,\partial N;\Z)$, the Thurston norm is
\[
 x_N(a)=\min_{[S]=a}\sum_j\max\{0,-\chi(S_j)\},
\]
where $S=\bigsqcup_jS_j$ ranges over properly embedded oriented surfaces representing $a$. Extending by homogeneity and continuity gives the Thurston seminorm on $H_2(N,\partial N;\R)$, equivalently on $H^1(N;\R)$ by Poincar\'e--Lefschetz duality. Its unit ball is
\[
 B_x(N)=\{a\in H_2(N,\partial N;\R):x_N(a)\le1\}.
\]
This is the usual unit-ball interpretation of the source's ``Thurston norm polytope.'' The integral homology lattice, modulo torsion, is part of the natural data.

\textbf{Problem (Agol).} Characterize the polytopes that occur as Thurston norm polytopes of finite-volume hyperbolic $3$-manifolds. In unit-ball terms, seek necessary and sufficient realizability conditions for a polytope $P$ in a real vector space with a specified integral lattice to be linearly identified, respecting the lattice, with $B_x(N)$ for some such $M$ and compact core $N$. This describes a natural full geometric realization problem; the excerpt does not explicitly specify whether equivalence is integral linear equivalence, real linear equivalence, or merely combinatorial equivalence, nor whether the dual norm polytope is intended instead. These conventions require confirmation. The standard definition above also uses orientability, which the extracted statement does not expressly impose.

The source reports that Thurston completed the rank-two case and described the higher-rank case as completely open in 2015. Here rank refers to the dimension of the norm's underlying homology/cohomology space under the usual interpretation. The supplied summary describes later computational and face-structure progress, without claiming a complete general realization characterization.
\subsection{Short English statement}
Determine exactly which polytopes arise from the Thurston norm of a finite-volume hyperbolic $3$-manifold, including the higher-rank cases. Specify whether the full lattice geometry, only combinatorial type, or a dual-polytope convention is being classified.
\subsection{Sources}
\begin{itemize}
\item Local UnsolvedMath registry, permanent ID 10900028, source code catalogue entry; attributed to Agol, no proposal year supplied.
\item Delp, Hoffoss and Manning, \emph{Problems In Groups, Geometry, and Three-Manifolds} (2015), subsection 5.3 (Agol), PDF page 4; \url{https://arxiv.org/pdf/1512.04620#page=4} and \url{https://arxiv.org/abs/1512.04620}. The registry records source-TeX extraction checked against rendered PDF text and access on 2026-07-29.
\item The source attributes the rank-two classification to Thurston; the research summary refers to his norm survey without full bibliographic data.
\item The summary mentions Dunfield--Kalelkar on hyperbolic homology classes, ``A cryptographic application of the Thurston norm'' with \url{https://arxiv.org/abs/1908.03504}, and work on norm faces via laminations with \url{https://arxiv.org/abs/2303.17665}. These are supplied references only; title/identifier matches, attribution, and relevance have not been checked here.
\end{itemize}
\end{document}
